\chapter{Cusp Forms and Representations of the Adele Group of \texorpdfstring{$\GL(2)$}{GL(2)}}

\label{ch:5}

Let $\bR(g)$ denote the natural representation of $\GL(2,\A)$ in the Hilbert space $\bL^2(\bG_\Q \setminus \bG_\A)$. This representation decomposes as a direct integral

\begin{equation}\label{eq:5.1}
\int \bL^2(\bG_\Q \setminus \bG_\A, \psi) \dd\psi\tag{5.1}
\end{equation}

according to the characters of the center of $\bG_\A$ trivial on $\Q^{\times}$. The resulting representation on each $\bL^2(\bG_\Q \setminus \bG_\A, \psi)$ will be denoted by $\bR^{\psi}(g)$.

In the language of representation theory one of the fundamental problems of the modern theory of automorphic forms is the following:

\textbf{PROBLEM:} Decompose $\bR^{\psi}(g)$ explicitly as a direct sum and integral of irreducible representations.

The rest of these Notes will be devoted in one form or another to a discussion of what is known concerning this decomposition.

It will turn out that the continuous spectrum of $\bR^{\psi}(g)$ is completely known, thanks to the theory of Eisenstein series developed by Selberg and Langlands. The discrete spectrum, on the other hand, is not. What is known is that this part of the spectrum is essentially exhausted by the space of cusp forms $\bL_0^2(\bG_\Q \setminus \bG_\A, \psi)$. Therefore, as in Section~\ref{ch:2}, it is natural to consider the restriction of $\bR^{\psi}(g)$ to $\bL_0^2(\bG_\Q \setminus \bG_\A, \psi)$ and to denote this restriction by $\bR_0^{\psi}(g)$.

In Section~\ref{ch:3} it was found that the classical cusp forms define special elements of $\bL_0^2(\bG_\Q \setminus \bG_\A, \psi)$. The purpose of the present section is to explain how these forms (and even the real analytic wave forms of Maass) correspond to irreducible constituents of $\bR_0^{\psi}$ in a natural and precise way.

Consequently classical cusp forms are nothing more nor less than special representations of $\GL(2,\A)$.

We shall also explain how the decomposition of $\bR_0^{\psi}(g)$ relates to the decomposition of the natural representation of $\SL(2,\R)$ in $\bL_0^2(\Gamma_0(N) \setminus \SL(2,\R))$. (See ``Notes and References.'')

\section{Preliminary Results on the Decomposition of \texorpdfstring{$\bR_0^{\psi}(g)$}{R0-psi(g)}}\label{sec:5a}

As already suggested, the main result here is the following adelic version of Part (ii) of Theorem~\ref{thm:2.6}.

\begin{theorem}\label{thm:5.1}

The representation $\bR_0^{\psi}(g)$ is the discrete direct sum of irreducible unitary representations each occurring with finite multiplicity.

\end{theorem}

This result is entirely functional analytic and a direct consequence of the spectral theory for compact self-adjoint operators. More precisely, its proof follows from the following two well-known lemmas.

\begin{lemma}\label{lem:5.2}

Suppose $f$ is any continuous compactly supported function on $\bG_\A$. Then the operator

\begin{equation}\label{eq:5.2}
\bR_0^{\psi}(f) = \int_{\bG_\A} f(g) \bR_0^{\psi}(g) \dd g\tag{5.2}
\end{equation}

is compact as an operator on $\bL_0^2(\bG_\Q \setminus \bG_\A, \psi)$.

\end{lemma}

\begin{proof}[Proof (Sketch)]

By definition,

\[
\bR_0^{\psi}(f)\phi(x) = \int_{\bG_\A} f(x^{-1}g)\phi(g) \dd g
\]

for each $\phi \in \bL_0^2(\bG_\Q \setminus \bG_\A, \psi)$. Therefore

\[
\begin{aligned}
\bR_0^{\psi}(f)\phi(x) &= \int_{N_\Q \setminus \bG_\A} \phi(g) \sum_{N_\Q} f(x^{-1}ng) \dd g\\
&= \int_{N_\Q \setminus \bG_\A} \phi(g) \bK(f, x, g) \dd g
\end{aligned}
\]

where

\begin{equation}\label{eq:5.3}
\bK(f,x,g) = \sum_{n \in N_\Q} f(x^{-1}ng) - \int_{N_\A} f(x^{-1}ng) \dd n\tag{5.3}
\end{equation}

since the cuspidal condition for $\phi$ implies that

\[
\int_{N_\Q \setminus \bG_\A} \phi(g) \int_{N_\A} f(x^{-1}ng) \dd n \dd g = 0.
\]

Now for each $c > 0$, let $S(c)$ denote the Siegel domain $\omega A_{\infty}^{+}(c)\bK$, where $\omega$ is a relatively compact subset of $N_\A A_\A$, and $A_{\infty}^{+}(c)$ is the set of matrices $\begin{pmatrix} a_1 & 0 \\ 0 & a_2 \end{pmatrix}$ in $A_{\infty}$ such that $\left|\frac{a_1}{a_2}\right| \geq c$. This domain may be chosen so that $\bG_\A = \bG_\Q S(c)$ and so that $\gamma S(c)$ intersects $S(c)$ for only finitely many $\gamma \in \bG_\Q$.

The burden of the proof of this lemma is to establish an estimate for the kernel $\bK(f, x, g)$ which implies that for any $x = u \begin{pmatrix} a_1 & 0 \\ 0 & a_2 \end{pmatrix} k$ in $S(c)$, and $\phi$ in $\bL_0^2(\bG_\Q \setminus \bG_\A, \psi)$,

\begin{equation}\label{eq:5.4}
|\bR_0^{\psi}(f)\phi(x)| \leq C_N \left|\frac{a_1}{a_2}\right|^{-N} \|\phi\|_2\tag{5.4}
\end{equation}

for every $N$. A complete proof of \eqref{eq:5.4} may be found in \cite{Langlands} or \cite{Godement2}. From \eqref{eq:5.4} it follows that the set

\[
\left\{\bR_0^{\psi}(f)\phi : \|\phi\|_2 \leq 1 \text{ in } \bL_0^2(\bG_\Q \setminus \bG_\A, \psi)\right\}
\]

is pre-compact. Thus the Lemma follows.

\end{proof}

\begin{lemma}\label{lem:5.3}

Suppose $U(g)$ is a unitary representation of a locally compact group $G$ on a Hilbert space $\bH$ with the property that the convolution operator

\[
U(f) = \int_G f(g) U(g) \dd g
\]

is compact for every continuous compactly supported function $f$ on $G$. Then $\bH$ is the direct sum of countably many invariant subspaces which occur with finite multiplicity.

\end{lemma}

The proof of this well-known Lemma may be found in \cite{GelfandGraevPyatetskiiShapiro} or \cite{Langlands}. Note that if $f(g) \geq 0$, and $f(g^{-1}) = f(g)$ then $\pi(f)$ is a self-adjoint compact operator and consequently has countable discrete spectrum and eigenvalues of finite multiplicity.

\begin{remark}\label{rem:5.4}

Theorem~\ref{thm:5.1} is more or less a representation-theoretic statement of the fact that any ``classical'' space of cusp forms is finite dimensional. (By a ``classical'' space of forms we understand a subspace of functions in $\bL^2(\Gamma \setminus \SL(2,\R))$ satisfying a $\bK$-and $\mathfrak{z}$-finiteness condition.)

\end{remark}

A refinement of Theorem~\ref{thm:5.1} is provided by the ``multiplicity one'' result of Jacquet and Langlands (Theorem~\ref{thm:5.7} below). This result asserts that any given irreducible unitary representation of $\GL(2,\A)$ occurs at most once in $\bR_0^{\psi}$. To describe its proof we need to argue as follows.

Let $\mathcal{H}(\bG_\A)$ denote the Hecke group algebra of $\bG_\A$ defined in the obvious way as the restricted tensor product

\[
\bigotimes_p \mathcal{H}(\bG_p)
\]

of the local Hecke group algebras $\mathcal{H}(\bG_p)$. (Restricted with respect to the unit element $1_p$ of $\mathcal{H}(\bG_p)$, the characteristic function of the compact subgroup $\bK_p$.)

Suppose we are given, for each $p$, an irreducible admissible representation $\pi_p'$ of $\mathcal{H}(\bG_p)$ on a vector space $V_p$ which for almost every $p$ contains the identity representation of $\bK_p$. Then we can define an irreducible representation

\[
\pi' = \bigotimes_p \pi_p'
\]

of $\mathcal{H}(\bG_\A)$ on $V = \bigotimes_p V_p$ in complete analogy with the construction of Section~\ref{ch:4}.C.

In particular, suppose $\pi$ is an irreducible representation of $\bG_\A$ in the subspace $\bH$ of $\bL_0^2(\bG_\Q \setminus \bG_\A, \psi)$ and $\bigotimes_p \pi_p$ is the factorization of $\pi$. For each $p$, put

\[
V_p = \bH_p^0 = \bigoplus_{\substack{\sigma_p\\ \left(\begin{array}{c} \text{algebraic}\\ \text{sum} \end{array}\right)}} \bH_p(\sigma_p)
\]

if $\pi_p$ acts on $\bH_p$. (The summation extends over all irreducible admissible representations $\sigma_p$ of $\bK_p$.) Let $\pi_p'$ denote the corresponding representation of $\mathcal{H}(\bG_p)$ on $V_p$. Then define a representation $\pi' = \bigotimes_p \pi_p'$ of $\mathcal{H}(\bG_\A)$ on $V = \bigotimes_p V_p$ as above. Clearly the isomorphism of $\bH$ with $\bigotimes_p \bH_p$ (which depended on a choice of $\bK_p$-fixed vector $\xi_p^0$ for almost every $p$, cf.\ Section~\ref{ch:4}.C) induces a natural isomorphism between $V$ and the space $\bH^0$ consisting of $\bK$-finite functions in $\bH$. Thus we have:

\begin{lemma}\label{lem:5.5}

Every irreducible constituent $\pi$ of the unitary representation $\bR_0^{\psi}(g)$ determines an (algebraically) irreducible representation $\pi'$ of the group algebra $\mathcal{H}(\bG_\A)$ on a space of $\bK$-finite functions in $\bL_0^2(\bG_\Q \setminus \bG_\A, \psi)$.

\end{lemma}

\begin{lemma}\label{lem:5.6}

Suppose $\bH$ is an irreducible subspace of $\bL_0^2(\bG_\Q \setminus \bG_\A, \psi)$. Then functions in the dense subspace $\bH^0$ (consisting of $\bK$-finite vectors) are $C^{\infty}$ (as functions of the infinite component) and rapidly decreasing in the sense that they satisfy the inequalities

\begin{equation}\label{eq:5.5}
|\phi(x)| \leq C_N(\phi) \left|\frac{a_1}{a_2}\right|^{-N}\tag{5.5}
\end{equation}

for $x = u \begin{pmatrix} a_1 & 0 \\ 0 & a_2 \end{pmatrix} k$ in the Siegel domain $S(c)$.

\end{lemma}

\begin{proof}

Since $\phi$ is $\bK$-finite, we may assume that $\phi$ belongs to some subspace $\bH(\sigma)$ of $\bH$ consisting of functions which transform under $\bK$ according to the irreducible representation $\sigma$. But then from the finite dimensionality of $\bH(\sigma)$ (remember admissibility!) it follows that

\[
\phi = \bR_0^{\psi}(f)\phi
\]

for some nice compactly supported function $f$ on $\bG_\A$ (which we may assume $C^{\infty}$). Therefore the first assertion of the Lemma is obvious and the second is a consequence of the estimate \eqref{eq:5.4}.

\end{proof}

\begin{theorem}\label{thm:5.7}

The multiplicity of an irreducible constituent $\pi$ of $\bR_0^{\psi}$ is one.

\end{theorem}

\begin{proof}[Proof (Sketch)]

Suppose $\bH$ is the space of $\pi$ in $\bL_0^2(\bG_\Q \setminus \bG_\A, \psi)$. Let $\pi'$ be the corresponding representation of $\mathcal{H}(\bG_\A)$ on $\bH^0$ constructed in Lemma~\ref{lem:5.5}. The functions $\phi$ in $\bH^0$ being nice by Lemma~\ref{lem:5.6}, they have Fourier expansions

\begin{equation}\label{eq:5.6}
\phi\!\left(\begin{pmatrix} 1 & x \\ 0 & 1 \end{pmatrix} g\right) = \sum_{\xi} \phi_{\xi}(g) \tau(\xi x)\tag{5.6}
\end{equation}

by Remark~\ref{rem:3.5}. But $\phi_{\xi}(g) = 0$ if $\xi = 0$, since $\pi$ occurs in the space of cusp forms.

On the other hand, if $\xi \neq 0$, it follows immediately from the left invariance of $\phi$ under $\bG_\Q$, and the invariance of $\dd x$ under $x \to \xi x$, that

\[
\begin{aligned}
\phi_{\xi}(g) &= \int_{\Q \setminus \A} \phi\!\left(\begin{pmatrix} 1 & x \\ 0 & 1 \end{pmatrix} g\right) \overline{\tau(\xi x)} \dd x\\
&= \int \phi\!\left(\begin{pmatrix} \xi & \xi x \\ 0 & 1 \end{pmatrix} g\right) \overline{\tau(\xi x)} \dd x\\
&= \int \phi\!\left(\begin{pmatrix} \xi & x \\ 0 & 1 \end{pmatrix} g\right) \overline{\tau(x)} \dd x\\
&= \int \phi\!\left(\begin{pmatrix} 1 & x \\ 0 & 1 \end{pmatrix} \begin{pmatrix} \xi & 0 \\ 0 & 1 \end{pmatrix} g\right) \overline{\tau(x)} \dd x
\end{aligned}
\]

therefore

\begin{equation}\label{eq:5.7}
\phi\!\left(\begin{pmatrix} 1 & x \\ 0 & 1 \end{pmatrix} g\right) = \sum_{\xi \neq 0} W_{\phi}\!\left(\begin{pmatrix} \xi & 0 \\ 0 & 1 \end{pmatrix} g\right) \tau(\xi x)\tag{5.7}
\end{equation}

where $W_{\phi}(g)$ is the $\bK$-finite $C^{\infty}$ function defined on $\bG_\A$ by

\begin{equation}\label{eq:5.8}
W_{\phi}(g) = \phi_1(g) = \int_{\Q \setminus \A} \phi\!\left(\begin{pmatrix} 1 & x \\ 0 & 1 \end{pmatrix} g\right) \overline{\tau(x)} \dd x.\tag{5.8}
\end{equation}

Now clearly the correspondence $\phi \to W_{\phi}$ maps $\bH^0$ one-to-one onto a space of functions $W(g)$ on $\bG_\A$ satisfying the following properties:

\begin{equation}\tag{5.9}\label{eq:5.9}
\begin{minipage}[t]{0.9\linewidth}
\begin{enumerate}[(i)]
\item $W\!\left(\begin{pmatrix} 1 & x \\ 0 & 1 \end{pmatrix} g\right) = \tau(x)W(g)$ for all $x \in \A$;
\item $W\!\left(\begin{pmatrix} x & 0 \\ 0 & x \end{pmatrix} g\right) = \psi(x)W(g)$ for all $x \in \A^{\times}$;
\item $W$ is right $\bK$-finite, $C^{\infty}$ as a function of the archimedean component of $\bG_\A$, and rapidly decreasing in the sense of \eqref{eq:5.5}; and
\item The space of these $W(g)$ is preserved by right convolution by elements of $\mathcal{H}(\bG_\A)$, and the resulting representation of $\mathcal{H}(\bG_\A)$ on this space is equivalent to $\pi'$.
\end{enumerate}
\end{minipage}
\end{equation}

The proof of Theorem~\ref{thm:5.7} then rests on the assertion that, given an irreducible representation $\pi'$, the space of functions $W(g)$ satisfying (i)--(iv) of \eqref{eq:5.9} is unique. This space of functions is called the Whittaker space (or Whittaker model) of $\pi'$. The fact that it is unique could be established now by ``elementary'' means. However, we prefer to postpone the proof until Section~\ref{ch:6} where there is further motivation for considering such models.

In any event (modulo this uniqueness result) the proof of Theorem~\ref{thm:5.7} is complete. Indeed if $\pi$ were to occur twice in $\bR_0^{\psi}$ we would have two ``Whittaker models'' for $\pi'$, a contradiction.

\end{proof}

\begin{remark}\label{rem:5.8}

Implicit throughout this discussion (and, in particular, the proof of Lemma~\ref{lem:5.5}) is the fact that any irreducible constituent of the unitary representation $\bR_0^{\psi}$ is admissible, hence factorizable. In \cite{GelfandGraevPyatetskiiShapiro} it seems to have been proved directly that every irreducible unitary representation of $\GL(2,\A)$ is factorizable (hence automatically admissible!).

\end{remark}

\section{Cusp Forms and Hecke Operators Revisited}\label{sec:5b}

In Section~\ref{ch:1} we described some special results of Hecke and Atkin-Lehner giving necessary and sufficient conditions for a cusp form to be completely determined by almost all its Fourier coefficients. These results may be generalized to cusp forms in $S_k(N, \psi)$ with arbitrary character $\psi$ and in this generality may be reformulated and strengthened as follows.

Let $S_k^-(N,\psi)$ denote the subspace of $S_k(N,\psi)$ generated by functions of the form $f(dz)$ where $f(z)$ belongs to $S_k(m,\psi)$ with $m$ a proper divisor of $N$ and $d$ any divisor of $N/m$. (This subspace is the zero subspace when $N = 1$, or more generally, when $\psi$ is primitive modulo $N$.) Let $S_k^+(N,\psi)$ denote the orthocomplement of $S_k^-(N,\psi)$ with respect to the Petersson inner product in $S_k(N,\psi)$. (This is the space of new forms in $S_k(N,\psi)$.)

Finally let $L$ denote the algebra generated by the Hecke operators $T(p)$ and $T(p)^*$ where $T(p)^*$ denotes the adjoint of $T(p)$ in $S_k(N, \psi)$ with respect to the Petersson metric. (Recall that $T(p)$ is necessarily normal only if $(p, N) = 1$.)

\begin{theorem}\label{thm:5.9}

\begin{enumerate}[(a)]

\item $S_k^+(N,\psi)$ is the largest $L$ invariant subspace of $S_k(N,\psi)$ on which $L$ is commutative and semi-simple;

\item Let $L^+$ denote the restriction of $L$ to $S_k^+(N,\psi)$. Then the $\C$-rank of $L^+$ equals the $\C$-dimension of $S_k^+(N,\psi)$;

\item Suppose that $f_1 \in S_k^+(N, \psi)$, $f_2 \in S_k^+(M, \psi)$, $T(p) f_i = a_p^i f_i$, $i = 1, 2$, for all $p$,

\[
a_p^1 = a_p^2
\]

for almost all $p$. Then $N = M$, and $f_1$ and $f_2$ coincide up to a constant multiple.

\end{enumerate}

\end{theorem}

The proof of Part (c) involves a comparison of the Dirichlet series attached to $f_1$ and $f_2$ and represents a great triumph of Hecke's theory.

There is a natural adelic generalization of Theorem~\ref{thm:5.9} which includes the case of the real analytic wave forms of Maass and which we shall presently describe.

Fix $\lambda$ an eigenvalue for the Casimir operator $\Delta$. (Thus $\lambda = -\frac{k}{2}\left(\frac{k}{2} - 1\right)$ or $\frac{1 - s^2}{4}$, with $s$ purely imaginary or purely real and between $-1$ and $1$.) For each irreducible representation $\sigma$ of $\operatorname{SO}(2) \subset \bK_{\infty}$ let $\bH(\psi, \lambda, N, \sigma)$ denote the subspace of functions $\phi$ in $\bL_0^2(\bG_\Q \setminus \bG_\A, \psi)$ such that

\begin{equation}\label{eq:5.10}
\Delta\phi = \lambda\phi\tag{5.10}
\end{equation}

and

\begin{equation}\label{eq:5.11}
\phi(gr(\theta)k_0) = \sigma(r(\theta))\psi(k_0)\phi(g)\tag{5.11}
\end{equation}

for all $g \in \bG_\A$, $r(\theta) \in \operatorname{SO}(2,\R)$, and $k_0 \in \bK_0^N$. (Here $\bK_0^N$ and $\psi$ are defined as in Section~\ref{ch:3}.A.) This subspace consists of cusp forms of type $(\psi, \lambda, N, \sigma)$.

\begin{example}\label{exa:5.10}

Suppose $\psi$ corresponds to a character of $(\Z/N\Z)^{\times}$, $\lambda = -\frac{k}{2}\left(\frac{k}{2} - 1\right)$, and $\sigma(r(\theta)) = e^{-ik\theta}$. Then $\bH(\psi, \lambda, N, \sigma)$ is isomorphic to the classical space $S_k(N, \psi)$ via the map $f \to \phi_f$ defined by \eqref{eq:3.4}.

\end{example}

\begin{example}\label{exa:5.11}

Suppose $\psi$ is trivial, $\lambda = \frac{1 - s^2}{4}$, and $\sigma$ is the trivial character of $\operatorname{SO}(2)$. Then $\bH(\psi, \lambda, N, \sigma)$ is isomorphic to the space $W_s(\Gamma_0(N))$ consisting of real-analytic cusp forms of level $N$. (The correspondence is simply $f(z) \leftrightarrow \phi(g) = f(g_\infty(i))$.)

\end{example}

Clearly $\bH(\psi,\lambda,N,\sigma)$ provides a meaningful generalization of the classical notion of cusp form. The natural analogue of the Petersson inner product for $\bH(\psi,\lambda,N,\sigma)$ is simply

\begin{equation}\label{eq:5.12}
(\phi_1, \phi_2) = \int_{Z_\A \bG_\Q \setminus \bG_\A} \phi_1(g) \overline{\phi_2(g)} \dd g\tag{5.12}
\end{equation}

and equipped with this inner product $\bH(\psi, \lambda, N, \sigma)$ is a finite dimensional Hilbert space.

Note that the integral in \eqref{eq:5.12} reduces to the classical integral in \eqref{eq:1.11} in case the hypotheses of Example~\ref{exa:5.10} are satisfied. Indeed in this case $\phi_1\overline{\phi_2}$ is a function on

\begin{equation}\label{eq:5.13}
Z_\A \bG_\Q \setminus \bG_\A / \bK_{\infty} \bK_0^N \cong \Gamma_0(N) \setminus \SL(2,\R) / \operatorname{SO}(2,\R)\tag{5.13}
\end{equation}

and $\phi_1\overline{\phi_2}(g) = f_1(z)\overline{f_2(z)}y^k$.

To state the analogue of Theorem~\ref{thm:5.9} for the spaces $\bH(\psi,\lambda,N,\sigma)$ we need first to introduce some Hecke operators appropriate for these spaces.

For each prime $p$, define $\widetilde{T}(p)$ and $\widetilde{T}(p)^{*}$ on $\bH(\psi, \lambda, \sigma, N)$ by

\begin{equation}\label{eq:5.14}
(\widetilde{T}(p)\phi)(g) = \int_{\bK_p} \psi^{-1}(k_p) \phi\!\left(gk_p \begin{pmatrix} p & 0 \\ 0 & 1 \end{pmatrix}\right) \dd k_p\tag{5.14}
\end{equation}

and

\begin{equation}\label{eq:5.15}
\widetilde{T}(p)^{*}\phi(g) = \int_{\bK_p} \psi^{-1}(k_p) \phi\!\left(gk_p \begin{pmatrix} p^{-1} & 0 \\ 0 & 1 \end{pmatrix}\right) \dd k_p\tag{5.15}
\end{equation}

where $\bK_p = \bG_p \cap \bK_0^N$. These operators are adjoint to one another with respect to the inner product defined by \eqref{eq:5.12}. If $\psi$ is trivial, $\widetilde{T}(p)$ coincides with the Hecke operator $\widetilde{T}(p)$ already introduced in Section~\ref{ch:3} (cf.\ \eqref{eq:3.15}).

Finally, let $\bH^-(\psi, \lambda, N, \sigma)$ denote the subspace of $\bH(\psi, \lambda, N, \sigma)$ generated by $\bH(\psi, \lambda, m, \sigma)$ for all proper divisors $m$ of $N$ and by elements of $\bH(\psi, \lambda, N, \sigma)$ which can be expressed as right transforms (by elements of $\bG_\A$) of functions in $\bH(\psi, \lambda, m, \sigma)$. (Caution: $\bH(\psi, \lambda, m, \sigma)$ is definitely not right $\bG_\A$-invariant!) Let $\bH^+(\psi, \lambda, N, \sigma)$ denote the orthocomplement of $\bH^-(\psi, \lambda, N, \sigma)$ in $\bH(\psi, \lambda, N, \sigma)$ with respect to the inner product \eqref{eq:5.12} and let $L$ denote the algebra over $\C$ generated by the operators $\widetilde{T}(p)$ and $\widetilde{T}(p)^{*}$ acting on $\bH(\psi, \lambda, N, \sigma)$.

\begin{theorem}\label{thm:5.12}

\begin{enumerate}[(a)]

\item $\bH^+(\psi,\lambda,N,\sigma)$ is the largest $L$-invariant subspace of $\bH(\psi,\lambda,N,\sigma)$ on which $L$ is commutative and semi-simple.

\item If $L^+$ denotes the restriction of $L$ to $\bH^+(\psi, \lambda, N, \sigma)$ then

\[
\rank_{\C}(L^+) = \dim_{\C} \bH^+(\psi, \lambda, N, \sigma).
\]

\item If $f_i \in \bH(\psi, \lambda, N_i, \sigma)$, and

\[
\widetilde{T}(p) f_i = a_p^i f_i \quad (i = 1, 2)
\]

for all $p$, then $f_1$ is a constant multiple of $f_2$ as soon as $a_p^1 = a_p^2$ for almost all $p$.

\end{enumerate}

\end{theorem}

\begin{remark}\label{rem:5.13}

The maximality assertion of Theorem~\ref{thm:5.9}(b) (and 5.12(b)) has the following significance. If a cusp form is a eigenfunction for all $\widetilde{T}(p)$ and $\widetilde{T}(p)^*$, then it is necessarily a new form.

\end{remark}

Theorem~\ref{thm:5.12} is proved in \cite{Miyake} more or less following the classical arguments of Hecke and Atkin-Lehner. The Corollary to this theorem which is of great interest to us and which represents a useful addition to the preliminary results on the decomposition of $\bR_0^{\psi}(g)$ is the following:

\begin{theorem}\label{thm:5.14}

Suppose $\pi^1$ and $\pi^2$ are two irreducible (unitary) constituents of $\bR_0^{\psi}$. Suppose also that the local components $\pi_p^1$ and $\pi_p^2$ are equivalent for all but finitely many finite places. Then $\pi^1 = \pi^2$.

\end{theorem}

\begin{proof}

Certainly it suffices to show that the spaces of $\pi_1'$ and $\pi_2'$ contain some function common to them both.

So suppose $\sigma$ is an irreducible constituent of the restriction of $\pi_{\infty}^1 = \pi_{\infty}^2$ to $\operatorname{SO}(2)$. Define the conductor of $\pi^i$ to be the integer

\begin{equation}\label{eq:5.16}
c(\pi^i) = \prod_{p < \infty} c(\pi_p^i)\tag{5.16}
\end{equation}

where $c(\pi_p^i)$ is the conductor of $\pi_p^i$ in the sense of Remark~\ref{rem:4.25}. Then let $\bH(\psi,\pi^i,c(\pi^i),\sigma)$ denote the subspace of functions in the space of $\pi^i$ in $\bL_0^2(\bG_\Q \setminus \bG_\A,\psi)$ which satisfy

\[
\phi(gr(\theta)k_0) = \sigma(r(\theta))\psi(k_0)\phi(g)
\]

for all $g \in \bG_\A$, $r(\theta) \in \operatorname{SO}(2,\R)$, and $k_0 \in \bK_0^{N_i}$ with $N_i = c(\pi^i)$. Since $\sigma$ occurs just once in $\pi_{\infty}^i$, it follows from Theorem~\ref{thm:4.24} that

\[
\dim \bH(\psi, \pi^i, c(\pi^i), \sigma) = 1
\]

for $i = 1$ and $2$. Consequently if $f_i$ spans $\bH(\psi, \pi^i, N_i, \sigma)$, then for some choice of scalars $\{a_p^i\}$,

\[
\widetilde{T}(p) f_i = a_p^i f_i
\]

for all $p$.

Now suppose the eigenvalue of the Casimir operator corresponding to $\pi_{\infty}^1 = \pi_{\infty}^2$ is $\lambda$ (cf.\ Lemma~\ref{lem:2.9}). Then clearly $\bH(\psi, \pi^i, N_i, \sigma) = \bH(\psi, \lambda, N_i, \sigma)$. Furthermore,

\[
a_p^1 = a_p^2
\]

for almost every $p$, by our hypothesis on $\pi_p^i$, the fact that almost every $\pi_p^i$ is class 1, and the identity \eqref{eq:4.16} which asserts that $a_p^i$ is completely determined by $\pi_p^i$ when $\pi_p^i$ is class 1. Therefore, by Part (c) of Theorem~\ref{thm:5.12}, $f_1$ is some multiple of $f_2$, and the proof of 5.14 is complete.

\end{proof}

\begin{remark}\label{rem:5.15}

A direct representation-theoretic proof of Theorem~\ref{thm:5.14} has been given by \cite{Casselman}. This proof depends upon Jacquet-Langlands' treatment of Dirichlet series associated to cusp forms and will be explained in Section~\ref{ch:6}.

\end{remark}

As a consequence of the proof of Theorem~\ref{thm:5.14} we can describe a representation-theoretic formulation of the Conjecture of Ramanujan-Petersson.

\begin{lemma}\label{lem:5.16}

Suppose $\pi$ is an irreducible unitary representation of $\bG_\A$ satisfying the following properties:

\begin{enumerate}[(1)]

\item $\pi = \bigotimes_p \pi_p$ occurs in $\bR_0^{\psi}$ with $\psi$ trivial on $\R_+^{\times}$;

\item the conductor of $\pi$ (in the sense of \eqref{eq:5.16}) is $N$;

\item $\pi_{\infty}$ is equivalent to the discrete series representation $\sigma(\mu_1, \mu_2)$ with $\mu_1\mu_2^{-1}(t) = t^{k-1}\sgn(t)$ (cf.\ Remark~\ref{rem:4.6}) and $\pi_p$ is equivalent to the class 1 representation $\pi(s_1^p, s_2^p)$ for all $p$ not dividing $N$.

\end{enumerate}

Let $f(z)$ be ``the'' new form in $S_k(N, \psi)$ with the property that $\phi_f(g)$ generates the (one-dimensional) subspace of functions $\phi$ in the space of $\pi$ such that

\begin{equation}\label{eq:5.17}
\phi(gr(\theta)k_0) = \sigma_k(r(\theta))\psi(k_0)\phi(g)\tag{5.17}
\end{equation}

for all $r(\theta) \in \operatorname{SO}(2)$ and $k_0 \in \bK_0^N$. THEN (at last)

\[
T(p)f = C_p f
\]

for all $p$, and for $p$ not dividing $N$,

\begin{equation}\label{eq:5.18}
C_p = p^{\frac{k-1}{2}}\left(p^{s_1} + p^{s_2}\right).\tag{5.18}
\end{equation}

\end{lemma}

\begin{proof}

The Lemma is an immediate consequence of Theorem~\ref{thm:4.23} (in particular \eqref{eq:4.16}) and Lemma~\ref{lem:3.7}. The one-dimensionality of the space of functions in $\bH$ satisfying \eqref{eq:5.17} is a consequence of Theorem~\ref{thm:4.24} as explained in the proof of Theorem~\ref{thm:5.14}.

\end{proof}

Recall that $T(p)$ is hermitian whenever $(p, N) = 1$ and $\psi$ is trivial. Therefore each ``eigenvalue'' $p^{\frac{k-1}{2}}(p^{s_1} + p^{s_2})$ should be real. But $\psi$ trivial implies $s_1 + s_2 = 0$. Consequently $p^{s_1}$ and $p^{s_2}$ are complex conjugates of one another when $\pi(s_1, s_2)$ belong to the continuous series.

When $\pi(s_1, s_2)$ belongs to the complementary series, both $p^{s_1}$ and $p^{s_2}$ are real. Therefore, in either case, $(p^{s_1} + p^{s_2})$ is indeed real.

\begin{proposition}\label{prop:5.17}

Suppose $\pi$ is an irreducible constituent of $\bR_0^{\psi}$ of the type described in Lemma~\ref{lem:5.16}. Then $\pi_p$ belongs to the continuous (as opposed to complementary) series of representations of $\bG_p$ for all $p$ not dividing $N$ if and only if the Ramanujan-Petersson Conjecture is true.

\end{proposition}

\begin{proof}

From \eqref{eq:5.18} it is clear that

\begin{equation}\label{eq:5.19}
|C_p| \leq 2p^{\frac{k-1}{2}} \text{ for } (p,N) = 1\tag{5.19}
\end{equation}

if and only if

\begin{equation}\label{eq:5.20}
|p^{s_1} + p^{s_2}| \leq 2.\tag{5.20}
\end{equation}

But \eqref{eq:5.19} is precisely the R-P Conjecture, and \eqref{eq:5.20} obtains iff $s_1$ and $s_2$ are pure imaginary and index a member of the continuous series. (Cf.\ Theorem~\ref{thm:4.27}: if $s_1, s_2$ indexes a complementary series representation, then $s_2 = s_1 - \sigma$ and $\operatorname{Re}(s_1) = \sigma/2$ for some $-1 < \sigma < 1$, $\sigma \neq 0$, so $|p^{s_1} + p^{s_2}| = |p^{s_1}(1 + p^{-\sigma})| = p^{\sigma/2} + p^{-\sigma/2} > 2$.)

\end{proof}

\begin{remark}\label{rem:5.18}

The Ramanujan-Petersson Conjecture has an obvious generalization to non-holomorphic cusp forms and for this conjecture Proposition~\ref{prop:5.17} (as well as Lemma~\ref{lem:5.16}) is still valid (with $\pi_{\infty}$ equivalent now to a principal series representation $\pi(\mu_1, \mu_2)$).

\end{remark}

\section{Some Explicit Features of the Correspondence Between Cusp Forms and Representations}\label{sec:5c}

Implicit in Subsection 5.B. is the fact that each irreducible constituent $\pi$ of $\bR_0^{\psi}$ uniquely determines a new form in $S_k(c(\pi),\psi)$. To complete the description of this correspondence between cusp forms in $S_k(N,\psi)$ and constituents of $\bR_0^{\psi}$ it is necessary to proceed in the reverse direction as well, namely from (not necessarily new) forms in $S_k(N,\psi)$ to representations $\pi$.

So suppose that $f(z)$ is an arbitrary form in $S_k(N, \psi)$. By Section~\ref{ch:1}.C. we may assume, without loss of generality, that $f$ is an eigenfunction for every $T(p)$ with $(p, N) = 1$. Say $T(p)f = a_p f$. Then the corresponding function

\[
\phi_f(g) = f(g_{\infty}(i)) j(g_{\infty}, i)^{-k} \psi(k_0)
\]

in $\bL_0^2(\bG_\Q \setminus \bG_\A, \psi)$ is an eigenfunction for $\widetilde{T}(p)$ for each $(p, N) = 1$. Now consider the subspace $\bH(f)$ of $\bL_0^2(\bG_\Q \setminus \bG_\A, \psi)$ spanned by translates of $\phi_f(g)$ under $\bG_\A$. Let $\pi_f$ denote the resulting representation of $\bG_\A$ in $\bH(f)$. Then $\pi_f$ is a direct sum of irreducible unitary representations of $\bG_\A$ each of which occurs in $\bR_0^{\psi}$ and our claim is that each of these irreducible representations has the same local component $\pi_p$ for $p = \infty$ and all finite $p$ not dividing $N$. From this it follows (by Theorem~\ref{thm:5.14}) that $\pi_f$ is the direct sum of mutually equivalent constituents of $\bR_0^{\psi}$ and hence (by ``multiplicity one'') that $\pi_f$ itself is an irreducible constituent of $\bR_0^{\psi}$.

To prove our claim we note first every right translate of $\phi_f(g)$ shares the same eigenvalue for ``the'' Casimir operator of $\bG_{\infty}$. Therefore the infinite component of each summand of $\pi_f$ is completely determined. More precisely, this component must be $\sigma(p, t)$ with $p = k-1$ and $t = 0$. (Compare the notation of Remark~\ref{rem:4.6}.)

Next we note that the space of each summand of $\pi_f$ contains a function which is right $\bK_p^N$ invariant and an eigenfunction of $\widetilde{T}(p)$ (with eigenvalue $a_p$) for $(p, N) = 1$. The $p$-component of each summand of $\pi_f$ is then the unique class 1 representation $\pi_p(\mu_1, \mu_2)$ of $\bG_p$ satisfying $\mu_1\mu_2 = \psi_p$ and $p^{\frac{k-1}{2}}(p^{s_1} + p^{s_2}) = a_p$ if $\mu_i(t) = |t|^{s_i}$. Thus our claim that the $p$-components of these summands are identical for almost every finite $p$ is immediate.

The results above, together with those of Subsection B, can be summarized as follows:

\begin{theorem}\label{thm:5.19}

\begin{enumerate}[(a)]

\item Suppose $f$ in $S_k(N, \psi)$ is an eigenfunction of $T(p)$ for every $p$ not dividing $N$. Define $\pi_f$ to be the representation of $\bG_\A$ determined by the translates of $\phi_f(g)$ (cf.\ 3.4). Then $\pi_f$ is irreducible.

\item Conversely, suppose $\pi$ is an irreducible constituent of $\bR_0^{\psi}$ with $\psi$ trivial on $(\R_+)^{\times}$. Define $\{f_{\pi}\}$ to be the equivalence class of the functions $f_{\pi}$ in $S_k(N,\psi)$ which share the same eigenvalue for every $T(p)$ with $p$ relatively prime to $N$, the conductor of $\pi$, and which are such that $\phi_{f_{\pi}}(g)$ is contained in the space of $\pi$. Then the resulting correspondence

\[
\{f\} \leftrightarrow \pi_f,
\]

between equivalence classes of cusp forms in the above sense, and irreducible constituents of $\bR_0^{\psi}$, is one-to-one.

\item The mapping $f \to \pi_f$, from functions in $S_k(N, \psi)$ to representations of $\bG_\A$, is one-to-one when restricted to new-forms in $S_k(N, \psi)$; it is not one-to-one when restricted to elements of any old-class in $S_k(N, \psi)$. (Cf.\ Section~\ref{ch:1}.D. for this terminology.)

\end{enumerate}

\end{theorem}

\begin{remark}\label{rem:5.20}

Theorem~\ref{thm:5.19} describes the nature of the correspondence between cusp forms and representations. In particular it describes the precise sense in which this correspondence is (and is not) one-to-one.

\end{remark}

Since the representation associated to a given form is uniquely determined by that form it is natural to ask if there is a recipe for describing the components of this representation in terms of the ``parameters'' of this form (its weight, level, character, and eigenvalues)? Unfortunately, the answer to this question, in general, seems to be no. However, in certain special cases such a recipe is possible; some of these cases are described below.

\begin{proposition}\label{prop:5.21}

Suppose $N = \prod p_i$ is a product of distinct primes. If $f \in S_k(\Gamma_0(N))$ is a new form and $T(p)f = a_p f$ for all $p$, then

\[
\pi_f = \bigotimes_p \pi_p
\]

can be described as follows:

\begin{enumerate}[(i)]

\item $\pi_{\infty}$ is the discrete series representation $\sigma(p, t)$ with $p = k-1$ and $t = 0$;

\item if $(p, N) = 1$, $\pi_p$ is the class 1 representation $\pi(\mu_1, \mu_2)$ such that $\mu_1\mu_2$ is trivial and $a_p = p^{\frac{k-1}{2}}(p^{s_1} + p^{s_2})$ if $\mu_i(t) = |t|^{s_i}$; and

\item if $p = p_i$, $\pi_p$ is ``the'' special representation of $\bG_p$ with trivial central character.

\end{enumerate}

\end{proposition}

\begin{proof}

Everything is clear from previous discussions except assertion (iii) which follows from the fact that $\phi_f$ is right invariant by

\[
\bK_p^1 = \left\{\begin{pmatrix} a & b \\ c & d \end{pmatrix} \in \GL(2, O_p) : c \in pO_p \right\}.
\]

That is, (iii) follows from the fact that $\pi_p$ has conductor 1 (or $pO_p$). Indeed it is clear from table \eqref{eq:4.20} that no supercuspidal representation has conductor 1 and that the principal series representation $\pi(\mu_1, \mu_2)$ cannot have conductor 1 since $\mu_1\mu_2$ must be trivial and therefore $\mu_1$ and $\mu_2$ are either both unramified or both ramified.

\end{proof}

\section*{Concluding Remarks}

\begin{enumerate}[(a)]

\item The correspondence between cusp forms and representations described in this last section is valid for non-holomorphic cusp forms as well. The only modifications required are the obvious ones ``at infinity.''

\item Everything in this section is valid in the setting of cusp forms defined over arbitrary global fields. We have worked exclusively within the context of $\Q$ only because this sometimes simplifies the notation and because it is from this classical context that we shall draw further examples in later sections.

\end{enumerate}

\section*{Notes and References}

\addcontentsline{toc}{subsection}{Notes and References}

Parts of our discussion in B and C are adapted from \cite{Casselman} and \cite{Miyake}.

The relation between the decomposition of $\bL_0^2(\bG_\Q \setminus \bG_\A, \psi)$ and $\bL_0^2(\Gamma_0(N) \setminus \SL(2,\R))$ can be explained as follows. Suppose $\pi$ is an irreducible constituent of $\bL_0^2(\bG_\Q \setminus \bG_\A, \psi)$ with conductor $c(\pi) = N$.

Then let $\bH_{\infty}$ denote the subspace of the space of $\pi$ consisting of functions in $\bL_0^2(\bG_\Q \setminus \bG_\A, \psi)$ satisfying

\[
\phi(gk_0) = \psi(k_0)\phi(g)
\]

for all $k_0 \in \bK_0^N$. ($\bK_0^N$ is defined by \eqref{eq:4.19}.) By Theorem~\ref{thm:4.24} this space is equivalent, as a $\GL(2,\R)$-module, to the representation $\pi_{\infty}$, the infinite component of $\pi$. But now

\[
Z_{\infty}^+ \bG_\Q \setminus \bG_\A / \bK_0^N \cong \Gamma_0'(N) \setminus \SL(2,\R)
\]

where $Z_{\infty}^+ = \left\{\begin{pmatrix} r & 0 \\ 0 & r \end{pmatrix} : r \in (\R_+)^{\times}\right\}$. Therefore, supposing only that $\psi$ is trivial on $(\R_+)^{\times}$, we conclude that $\bH_{\infty}$, as a $\SL(2,\R)$-module, is isomorphic to an (essentially) irreducible subspace of $\bL_0^2(\Gamma_0(N) \setminus \SL(2,\R))$ which is equivalent to the restriction of $\pi_{\infty}$ to $\SL(2,\R)$. We say ``essentially'' irreducible because the restriction of $\pi_{\infty}$ to $\SL(2,\R)$ decomposes into two inequivalent pieces if (and only if) $\pi_{\infty}$ is a discrete series representation. (For details, see, for example \cite{Gelbart} Chapter II, Section~\ref{ch:1}.) In any event (either piece of) the restriction of $\pi_{\infty}$ to $\SL(2,\R)$ is an irreducible constituent of $\bL_0^2(\Gamma_0(N) \setminus \SL(2,\R))$.

In \cite{Langlands} and \cite{Godement2} Theorem~\ref{thm:5.1} is proved in the generality of an arbitrary reductive group $G$. The fact that $\bL_0^2(\bG_\Q \setminus \bG_\A, \psi)$ essentially exhausts the discrete spectrum follows from the theory of Eisenstein series and is discussed in Section~\ref{ch:8} of these Notes. The characterization of those $\pi$ which occur in $\bL_0^2(\bG_\Q \setminus \bG_\A, \psi)$ constitutes Jacquet-Langlands' reinterpretation of Hecke's theory and is the proper subject matter of Section~\ref{ch:6}.

Another group-theoretic treatment of the Ramanujan-Petersson Conjecture (Theorem) appears in \cite{Satake}.
